% \documentclass{article}
% \usepackage{graphicx} % Required for inserting images

% \title{bci}
% \author{INTERFACEs VISTEC}
% \date{September 2026}

% \begin{document}

% \maketitle

% \section{Introduction}

% \end{document}
% Options for packages loaded elsewhere
\PassOptionsToPackage{unicode}{hyperref}
\PassOptionsToPackage{hyphens}{url}
\PassOptionsToPackage{dvipsnames,svgnames,x11names}{xcolor}
\documentclass[
  10pt,
]{article}
\usepackage{xcolor}
\usepackage[margin=1in]{geometry}
\usepackage{amsmath,amssymb}
\setcounter{secnumdepth}{-\maxdimen} % remove section numbering
\usepackage{iftex}
\ifPDFTeX
  \usepackage[T1]{fontenc}
  \usepackage[utf8]{inputenc}
  \usepackage{textcomp} % provide euro and other symbols
\else % if luatex or xetex
  \usepackage{unicode-math} % this also loads fontspec
  \defaultfontfeatures{Scale=MatchLowercase}
  \defaultfontfeatures[\rmfamily]{Ligatures=TeX,Scale=1}
\fi
\usepackage{lmodern}
\ifPDFTeX\else
  % xetex/luatex font selection
\fi
% Use upquote if available, for straight quotes in verbatim environments
\IfFileExists{upquote.sty}{\usepackage{upquote}}{}
\IfFileExists{microtype.sty}{% use microtype if available
  \usepackage[]{microtype}
  \UseMicrotypeSet[protrusion]{basicmath} % disable protrusion for tt fonts
}{}
\makeatletter
\@ifundefined{KOMAClassName}{% if non-KOMA class
  \IfFileExists{parskip.sty}{%
    \usepackage{parskip}
  }{% else
    \setlength{\parindent}{0pt}
    \setlength{\parskip}{6pt plus 2pt minus 1pt}}
}{% if KOMA class
  \KOMAoptions{parskip=half}}
\makeatother
\usepackage{color}
\usepackage{fancyvrb}
\newcommand{\VerbBar}{|}
\newcommand{\VERB}{\Verb[commandchars=\\\{\}]}
\DefineVerbatimEnvironment{Highlighting}{Verbatim}{commandchars=\\\{\}}
% Add ',fontsize=\small' for more characters per line
\usepackage{framed}
\definecolor{shadecolor}{RGB}{248,248,248}
\newenvironment{Shaded}{\begin{snugshade}}{\end{snugshade}}
\newcommand{\AlertTok}[1]{\textcolor[rgb]{0.94,0.16,0.16}{#1}}
\newcommand{\AnnotationTok}[1]{\textcolor[rgb]{0.56,0.35,0.01}{\textbf{\textit{#1}}}}
\newcommand{\AttributeTok}[1]{\textcolor[rgb]{0.13,0.29,0.53}{#1}}
\newcommand{\BaseNTok}[1]{\textcolor[rgb]{0.00,0.00,0.81}{#1}}
\newcommand{\BuiltInTok}[1]{#1}
\newcommand{\CharTok}[1]{\textcolor[rgb]{0.31,0.60,0.02}{#1}}
\newcommand{\CommentTok}[1]{\textcolor[rgb]{0.56,0.35,0.01}{\textit{#1}}}
\newcommand{\CommentVarTok}[1]{\textcolor[rgb]{0.56,0.35,0.01}{\textbf{\textit{#1}}}}
\newcommand{\ConstantTok}[1]{\textcolor[rgb]{0.56,0.35,0.01}{#1}}
\newcommand{\ControlFlowTok}[1]{\textcolor[rgb]{0.13,0.29,0.53}{\textbf{#1}}}
\newcommand{\DataTypeTok}[1]{\textcolor[rgb]{0.13,0.29,0.53}{#1}}
\newcommand{\DecValTok}[1]{\textcolor[rgb]{0.00,0.00,0.81}{#1}}
\newcommand{\DocumentationTok}[1]{\textcolor[rgb]{0.56,0.35,0.01}{\textbf{\textit{#1}}}}
\newcommand{\ErrorTok}[1]{\textcolor[rgb]{0.64,0.00,0.00}{\textbf{#1}}}
\newcommand{\ExtensionTok}[1]{#1}
\newcommand{\FloatTok}[1]{\textcolor[rgb]{0.00,0.00,0.81}{#1}}
\newcommand{\FunctionTok}[1]{\textcolor[rgb]{0.13,0.29,0.53}{\textbf{#1}}}
\newcommand{\ImportTok}[1]{#1}
\newcommand{\InformationTok}[1]{\textcolor[rgb]{0.56,0.35,0.01}{\textbf{\textit{#1}}}}
\newcommand{\KeywordTok}[1]{\textcolor[rgb]{0.13,0.29,0.53}{\textbf{#1}}}
\newcommand{\NormalTok}[1]{#1}
\newcommand{\OperatorTok}[1]{\textcolor[rgb]{0.81,0.36,0.00}{\textbf{#1}}}
\newcommand{\OtherTok}[1]{\textcolor[rgb]{0.56,0.35,0.01}{#1}}
\newcommand{\PreprocessorTok}[1]{\textcolor[rgb]{0.56,0.35,0.01}{\textit{#1}}}
\newcommand{\RegionMarkerTok}[1]{#1}
\newcommand{\SpecialCharTok}[1]{\textcolor[rgb]{0.81,0.36,0.00}{\textbf{#1}}}
\newcommand{\SpecialStringTok}[1]{\textcolor[rgb]{0.31,0.60,0.02}{#1}}
\newcommand{\StringTok}[1]{\textcolor[rgb]{0.31,0.60,0.02}{#1}}
\newcommand{\VariableTok}[1]{\textcolor[rgb]{0.00,0.00,0.00}{#1}}
\newcommand{\VerbatimStringTok}[1]{\textcolor[rgb]{0.31,0.60,0.02}{#1}}
\newcommand{\WarningTok}[1]{\textcolor[rgb]{0.56,0.35,0.01}{\textbf{\textit{#1}}}}
\usepackage{longtable,booktabs,array}
\newcounter{none} % for unnumbered tables
\usepackage{calc} % for calculating minipage widths
% Correct order of tables after \paragraph or \subparagraph
\usepackage{etoolbox}
\makeatletter
\patchcmd\longtable{\par}{\if@noskipsec\mbox{}\fi\par}{}{}
\makeatother
% Allow footnotes in longtable head/foot
\IfFileExists{footnotehyper.sty}{\usepackage{footnotehyper}}{\usepackage{footnote}}
\makesavenoteenv{longtable}
\setlength{\emergencystretch}{3em} % prevent overfull lines
\providecommand{\tightlist}{%
  \setlength{\itemsep}{0pt}\setlength{\parskip}{0pt}}
% --- Unicode glyphs that pdfLaTeX does not know out of the box -------------
% XeLaTeX/LuaLaTeX get these from unicode-math; pdfLaTeX needs them declared.
\ifPDFTeX
  \usepackage{newunicodechar}
  \newunicodechar{→}{\ensuremath{\rightarrow}}
  \newunicodechar{∈}{\ensuremath{\in}}
\fi
\usepackage{bookmark}
\IfFileExists{xurl.sty}{\usepackage{xurl}}{} % add URL line breaks if available
\urlstyle{same}
\hypersetup{
  pdftitle={Homework --- Build and Benchmark Your Own MI-EEG Pipeline},
  colorlinks=true,
  linkcolor={blue},
  filecolor={Maroon},
  citecolor={Blue},
  urlcolor={blue},
  pdfcreator={LaTeX via pandoc}}

\title{Homework\\Build and Benchmark Your Own MI-EEG Pipeline}
\author{}
\date{}

\begin{document}
\maketitle

{
\setcounter{tocdepth}{2}
\tableofcontents
}

\begin{center}\rule{0.5\linewidth}{0.5pt}\end{center}

\textbf{Dataset:} BCI Competition IV-2a (BCIC2a) · \textbf{Setting:}
subject-dependent · \textbf{Task:} 2-class motor imagery (left hand vs
right hand)

You will design your own EEG preprocessing pipeline and your own neural
network, plug both into this repository, and benchmark them against
\textbf{MixNet} under exactly the same protocol --- same 9 subjects,
same 5 folds, same test session, same metrics.



\section{1. Learning objectives}\label{learning-objectives}

By the end of this assignment you should be able to:

\begin{enumerate}
\def\labelenumi{\arabic{enumi}.}
\tightlist
\item
  Implement a feature-extraction pipeline and a deep model against a
  fixed experimental API instead of a one-off notebook.
\item
  Run a controlled comparison against a strong published baseline and
  report it honestly, including where your method loses.
\item
  Analyse a result instead of only reporting it: read the per-subject
  table against the paired significance test, and form a hypothesis for
  the pattern you see.
\end{enumerate}



\section{2. The dataset}\label{the-dataset}

{\def\LTcaptype{none} % do not increment counter
\begin{longtable}[]{@{}
  >{\raggedright\arraybackslash}p{(\linewidth - 2\tabcolsep) * \real{0.3}}
  >{\raggedright\arraybackslash}p{(\linewidth - 2\tabcolsep) * \real{0.8}}@{}}
\toprule\noalign{}
\begin{minipage}[b]{\linewidth}\raggedright
Property
\end{minipage} & \begin{minipage}[b]{\linewidth}\raggedright
Description
\end{minipage} \\
\midrule\noalign{}
\endhead
\bottomrule\noalign{}
\endlastfoot
Subjects & 9 (\texttt{A01}--\texttt{A09}) \\
Sessions & \texttt{T} (training) and \texttt{E} (evaluation), recorded
on different days \\
Classes used & 2 --- left hand (\texttt{0}), right hand (\texttt{1}) \\
Trials & 144 per session per subject (72 per class) \\
Channels & 20 pre-selected motor-cortex channels (see
\texttt{mixnet/preprocessing/config.py}) \\
Original rate & 250 Hz, downsampled to 100 Hz by the provided loader \\
MI window & 4 S (2 s → 6 s after the cue) \\
\end{longtable}
}

\textbf{Subject-dependent} means each subject gets their own model. For
every subject, session \texttt{T} is split into train/validation with
5-fold stratified cross-validation, and session \texttt{E} is the
\textbf{untouched test set} of every fold. Never train on session
\texttt{E}.



\section{3. What is given vs what you
write}\label{what-is-given-vs-what-you-write}

{\def\LTcaptype{none} % do not increment counter
\begin{longtable}[]{@{}
  >{\raggedright\arraybackslash}p{(\linewidth - 2\tabcolsep) * \real{0.5000}}
  >{\raggedright\arraybackslash}p{(\linewidth - 2\tabcolsep) * \real{0.5000}}@{}}
\toprule\noalign{}
\begin{minipage}[b]{\linewidth}\raggedright
File
\end{minipage} & \begin{minipage}[b]{\linewidth}\raggedright
Description
\end{minipage} \\
\midrule\noalign{}
\endhead
\bottomrule\noalign{}
\endlastfoot
\texttt{mixnet/preprocessing/BCIC2a/HW\_prep.py} & \textbf{YOU WRITE}
--- 2 functions \\
\texttt{mixnet/models/HW\_Net.py} & \textbf{YOU WRITE} --- 2
functions \\
\texttt{experiments/configs/HW\_Net.py} & \textbf{YOU TUNE} --- shapes +
hyper-parameters \\
\texttt{experiments/prep\_HW.py} & given --- runs your pipeline (\textbf{NOTE:
add your own parameters}) \\
\texttt{experiments/run\_HW\_Net.py} & given --- 5-fold training/eval
driver \\
\texttt{experiments/benchmark.py} & given --- aggregation + statistics,
no edits needed \\
everything under \texttt{mixnet/} else & given --- do not modify \\
\end{longtable}
}

Each file you edit starts with a header block explaining the contract it
must satisfy. Read those headers first --- they answer most questions.



\section{4. Setup}\label{setup}

\begin{quote}
\textbf{Run every command below from the \texttt{experiments/}
directory} unless stated otherwise. The scripts resolve
\texttt{datasets/} and \texttt{logs/} relative to the current working
directory, and both are git-ignored.
\end{quote}

\subsection{4.1 Prerequisites}\label{prerequisites}

\begin{itemize}
\tightlist
\item
  Python \textbf{3.8.10} (the package declares
  \texttt{\textgreater{}=3.8,\ \textless{}=3.10.4})
\item
  An NVIDIA GPU with CUDA --- \texttt{mixnet/models/base.py} reads GPU
  memory during evaluation, so a GPU is effectively required
\item
  \texttt{git}
\end{itemize}

\subsection{4.2 Fork and clone}\label{fork-and-clone}

Fork \texttt{https://github.com/Max-Phairot-A/MixNet} to your own GitHub
account so you can commit your work, then:

\begin{Shaded}
\begin{Highlighting}[]
\FunctionTok{git}\NormalTok{ clone https://github.com/}\OperatorTok{\textless{}}\NormalTok{your{-}github{-}username}\OperatorTok{\textgreater{}}\NormalTok{/MixNet.git}
\BuiltInTok{cd}\NormalTok{ MixNet}
\FunctionTok{git}\NormalTok{ switch homework}
\end{Highlighting}
\end{Shaded}

Work on that branch.

\subsection{4.3 Create the environment}\label{create-the-environment}

The reference environment is the TensorFlow 2.7.0 GPU image:

\begin{Shaded}
\begin{Highlighting}[]
\ExtensionTok{docker}\NormalTok{ pull tensorflow/tensorflow:2.7.0{-}gpu}
\ExtensionTok{docker}\NormalTok{ run }\AttributeTok{{-}ti} \AttributeTok{{-}{-}gpus}\NormalTok{ all }\AttributeTok{{-}{-}name}\NormalTok{ mixnet\_container }\AttributeTok{{-}v} \VariableTok{$(}\BuiltInTok{pwd}\VariableTok{)}\NormalTok{:/workspace }\DataTypeTok{\textbackslash{}}
\NormalTok{    docker.io/tensorflow/tensorflow:2.7.0{-}gpu bash}
\BuiltInTok{cd}\NormalTok{ /workspace}
\end{Highlighting}
\end{Shaded}

Or with conda, if you already have a working CUDA setup:

\begin{Shaded}
\begin{Highlighting}[]
\ExtensionTok{conda}\NormalTok{ create }\AttributeTok{{-}n}\NormalTok{ mixnet python=3.8.10 }\AttributeTok{{-}y}
\ExtensionTok{conda}\NormalTok{ activate mixnet}
\ExtensionTok{pip}\NormalTok{ install tensorflow{-}gpu==2.7.0}
\end{Highlighting}
\end{Shaded}

\subsection{4.4 Install the package --- read this
carefully}\label{install-the-package-read-this-carefully}

\begin{Shaded}
\begin{Highlighting}[]
\ExtensionTok{pip}\NormalTok{ install }\AttributeTok{{-}r}\NormalTok{ requirements.txt}
\ExtensionTok{pip}\NormalTok{ install }\AttributeTok{{-}e}\NormalTok{ .         }
\end{Highlighting}
\end{Shaded}

\begin{quote}
% \textbf{Use pip install -e .}

\texttt{pip\ install\ -e\ .} installs in \emph{editable} mode: it links
to your working copy, so your edits take effect immediately. If you must
use \texttt{pip\ install\ .}, you have to re-run it after \textbf{every}
change under \texttt{mixnet/}.
\end{quote}

Verify the install points at your clone:

\begin{Shaded}
\begin{Highlighting}[]
\ExtensionTok{python} \AttributeTok{{-}c} \StringTok{"import mixnet; print(mixnet.\_\_file\_\_)"}
\CommentTok{\# should print \textless{}your clone\textgreater{}/mixnet/\_\_init\_\_.py — NOT a site{-}packages path}
\end{Highlighting}
\end{Shaded}

\subsection{4.5 Download the data}\label{download-the-data}

\begin{Shaded}
\begin{Highlighting}[]
\BuiltInTok{cd}\NormalTok{ experiments}
\ExtensionTok{python}\NormalTok{ download\_datasets.py }\AttributeTok{{-}{-}dataset} \StringTok{\textquotesingle{}BCIC2a\textquotesingle{}}
\end{Highlighting}
\end{Shaded}

This writes
\texttt{experiments/datasets/BCIC2a/raw/A0\{1..9\}\{T,E\}.mat}. If the
download fails, fetch the files manually from
\url{http://bnci-horizon-2020.eu/database/data-sets} (dataset 001-2014)
and drop them into that folder.



\section{5. Reproduce the MixNet baseline
first}\label{reproduce-the-mixnet-baseline-first}

Do this \textbf{before} writing any code. It confirms your environment
works and gives you the number you have to beat.

\begin{Shaded}
\begin{Highlighting}[]
\BuiltInTok{cd}\NormalTok{ experiments}

\CommentTok{\# MixNet\textquotesingle{}s own preprocessing (spectral{-}spatial signals), subject{-}dependent only}
\ExtensionTok{python}\NormalTok{ prep\_spectral\_spatial\_signals.py }\AttributeTok{{-}{-}dataset} \StringTok{\textquotesingle{}BCIC2a\textquotesingle{}} \AttributeTok{{-}{-}setting} \StringTok{\textquotesingle{}dependent\textquotesingle{}}

\CommentTok{\# MixNet with the paper\textquotesingle{}s optimal BCIC2a subject{-}dependent hyper{-}parameters}
\ExtensionTok{python}\NormalTok{ run\_MixNet.py }\AttributeTok{{-}{-}model\_name} \StringTok{\textquotesingle{}MixNet\textquotesingle{}} \AttributeTok{{-}{-}dataset} \StringTok{\textquotesingle{}BCIC2a\textquotesingle{}} \DataTypeTok{\textbackslash{}}
    \AttributeTok{{-}{-}train\_type} \StringTok{\textquotesingle{}subject\_dependent\textquotesingle{}} \AttributeTok{{-}{-}data\_type} \StringTok{\textquotesingle{}spectral\_spatial\_signals\textquotesingle{}} \DataTypeTok{\textbackslash{}}
    \AttributeTok{{-}{-}adaptive\_gradient}\NormalTok{ True }\AttributeTok{{-}{-}policy} \StringTok{\textquotesingle{}HistoricalTangentSlope\textquotesingle{}} \DataTypeTok{\textbackslash{}}
    \AttributeTok{{-}{-}log\_dir} \StringTok{\textquotesingle{}logs\textquotesingle{}} \AttributeTok{{-}{-}num\_class}\NormalTok{ 2 }\AttributeTok{{-}{-}GPU}\NormalTok{ 0 }\DataTypeTok{\textbackslash{}}
    \AttributeTok{{-}{-}margin}\NormalTok{ 1.0 }\AttributeTok{{-}{-}n\_component}\NormalTok{ 2 }\AttributeTok{{-}{-}warmup}\NormalTok{ 7}
\end{Highlighting}
\end{Shaded}

Record the resulting accuracy and F1 with:

\begin{Shaded}
\begin{Highlighting}[]
\ExtensionTok{python}\NormalTok{ benchmark.py}
\end{Highlighting}
\end{Shaded}

\textbf{Deliverable checkpoint:} paste the MixNet summary table into
your report. That is your baseline; every later claim is relative to it.



\section{6. Task 1 --- your preprocessing
pipeline}\label{task-1-your-preprocessing-pipeline}

\textbf{File:} \texttt{mixnet/preprocessing/BCIC2a/HW\_prep.py}

Implement two functions:

{\def\LTcaptype{none} % do not increment counter
\begin{longtable}[]{@{}
  >{\raggedright\arraybackslash}p{(\linewidth - 4\tabcolsep) * \real{0.3333}}
  >{\raggedright\arraybackslash}p{(\linewidth - 4\tabcolsep) * \real{0.3333}}
  >{\raggedright\arraybackslash}p{(\linewidth - 4\tabcolsep) * \real{0.3333}}@{}}
\toprule\noalign{}
\begin{minipage}[b]{\linewidth}\raggedright
Function
\end{minipage} & \begin{minipage}[b]{\linewidth}\raggedright
Called with
\end{minipage} & \begin{minipage}[b]{\linewidth}\raggedright
Must return
\end{minipage} \\
\midrule\noalign{}
\endhead
\bottomrule\noalign{}
\endlastfoot
\texttt{fit\_transform(X,\ y,\ pick\_smp\_freq,\ **kwargs)} & the
training fold only & \texttt{(X\_out,\ state)} \\
\texttt{transform(X,\ state,\ pick\_smp\_freq,\ **kwargs)} & validation
and test folds & \texttt{X\_out} \\
\end{longtable}
}

You receive \texttt{X} of shape \texttt{(n\_trials,\ 20,\ 400)} --- 20
channels, 4 s at 100 Hz --- and labels \texttt{y\ ∈\ \{0,\ 1\}}. You
return features of shape \texttt{(n\_trials,\ ...)}; the trailing shape
is yours to choose.

\subsection{The rule that carries the most
marks}\label{the-rule-that-carries-the-most-marks}

\textbf{Anything that learns from data must be fitted inside
\texttt{fit\_transform} on the training fold only, and merely applied
inside \texttt{transform}.} CSP filters, PCA, a mean/std scaler, a
whitening matrix, a channel-selection criterion --- all of these leak if
you fit them on the full dataset before splitting. Leakage inflates
accuracy and makes your comparison against MixNet meaningless.

Stateless operations (band-pass filtering, cropping, resampling,
log-variance) are safe to apply to each split independently.

\subsection{Approaches you may take}\label{approaches-you-may-take}

\begin{itemize}
\tightlist
\item
  filter bank + CSP (\texttt{from\ mixnet.preprocessing\ import\ FBCSP})
\item
  band-power / log-variance features per channel per band
\item
  Riemannian covariance + tangent-space projection
\item
  STFT or wavelet time-frequency maps
\item
  band-pass + per-channel z-scoring (a simple, respectable baseline)
\end{itemize}

Helpers already in the repo:
\texttt{mixnet.utils.butter\_bandpass\_filter}, \texttt{resampling},
\texttt{psd\_welch}; \texttt{mixnet.preprocessing.FBCSP},
\texttt{SpectralSpatialMapping}.

NOTE: You can also add your own helper functions in this file or in
\texttt{mixnet/preprocessing/BCIC2a/HW\_prep\_helpers.py} if you want to
keep this file clean.

\subsection{Run it}\label{run-it}

\begin{Shaded}
\begin{Highlighting}[]
\BuiltInTok{cd}\NormalTok{ experiments}
\ExtensionTok{python}\NormalTok{ prep\_HW.py}
\end{Highlighting}
\end{Shaded}

Output lands in
\texttt{experiments/datasets/BCIC2a/HW\_prep/2\_class/subject\_dependent/}.
Expect 6 \texttt{.npy} files per fold for a subject (30 files per
subject) --- 270 files in total for 9 subjects.



\section{7. Task 2 --- your model}\label{task-2-your-model}

\textbf{File:} \texttt{mixnet/models/HW\_Net.py}

Implement \texttt{\_config()} (hyper-parameters) and \texttt{build()}
(the architecture). The training loops (\texttt{train\_step} /
\texttt{val\_step} / \texttt{test\_step} / \texttt{pred\_step}) are
already written for you. Note that you may rewrite these steps if you go
multi-task or use other loss functions. See MixNet for reference.

\subsection{Contract}\label{contract}

\begin{itemize}
\tightlist
\item
  \textbf{Input:} one tensor of shape \texttt{self.input\_shape} (no
  batch axis)
\item
  \textbf{Output:} exactly \textbf{one} tensor of shape
  \texttt{(batch,\ num\_class)} ending in
  \texttt{layers.Activation(\textquotesingle{}softmax\textquotesingle{})}
  --- multi-output models take a different code path in
  \texttt{mixnet/models/base.py} and will not work
\item
  \textbf{Name:} \texttt{name=self.model\_name}; never
  \texttt{\textquotesingle{}MixNet\textquotesingle{}} or
  \texttt{\textquotesingle{}MIN2Net\textquotesingle{}}, because
  \texttt{base.py} branches on those names
\item
  Keep the \texttt{load\_weights} block, or \texttt{evaluate()} will
  silently score an untrained network
\end{itemize}

\subsection{Get the plumbing working
first}\label{get-the-plumbing-working-first}

Before designing anything clever, paste this into \texttt{build()} and
run one subject end to end:

\begin{Shaded}
\begin{Highlighting}[]
\NormalTok{input1  }\OperatorTok{=}\NormalTok{ layers.Input(shape}\OperatorTok{=}\VariableTok{self}\NormalTok{.input\_shape)}
\NormalTok{x       }\OperatorTok{=}\NormalTok{ layers.Flatten()(input1)}
\NormalTok{x       }\OperatorTok{=}\NormalTok{ layers.Dense(}\VariableTok{self}\NormalTok{.num\_class)(x)}
\NormalTok{softmax }\OperatorTok{=}\NormalTok{ layers.Activation(}\StringTok{\textquotesingle{}softmax\textquotesingle{}}\NormalTok{, name}\OperatorTok{=}\StringTok{\textquotesingle{}softmax\textquotesingle{}}\NormalTok{)(x)}
\NormalTok{model   }\OperatorTok{=}\NormalTok{ Model(inputs}\OperatorTok{=}\NormalTok{input1, outputs}\OperatorTok{=}\NormalTok{softmax, name}\OperatorTok{=}\VariableTok{self}\NormalTok{.model\_name)}
\end{Highlighting}
\end{Shaded}

Then replace it. \texttt{mixnet/models/EEGNet.py} and
\texttt{DeepConvNet.py} are two complete worked examples of the same
contract.

\subsection{If you go multi-task, the step functions are no longer
free}\label{if-you-go-multi-task-the-step-functions-are-no-longer-free}

The \texttt{train\_step} / \texttt{val\_step} / \texttt{test\_step} /
\texttt{pred\_step} methods given to you are written for a
\textbf{single-task} model: one output, one loss (cross-entropy). MixNet
is multi-task --- it optimises reconstruction (MSE) + triplet +
cross-entropy together with adaptive loss weights --- and you may do the
same, but then \textbf{you have to rewrite all four steps}, plus three
other things that must change with them:

{\def\LTcaptype{none} % do not increment counter
\begin{longtable}[]{@{}
  >{\raggedright\arraybackslash}p{(\linewidth - 2\tabcolsep) * \real{0.5000}}
  >{\raggedright\arraybackslash}p{(\linewidth - 2\tabcolsep) * \real{0.5000}}@{}}
\toprule\noalign{}
\begin{minipage}[b]{\linewidth}\raggedright
What
\end{minipage} & \begin{minipage}[b]{\linewidth}\raggedright
Change
\end{minipage} \\
\midrule\noalign{}
\endhead
\bottomrule\noalign{}
\endlastfoot
\texttt{build()} & return a \textbf{list} of outputs, classifier softmax
\textbf{last} \\
\texttt{configs/HW\_Net.py} & one entry per task, same order, in
\texttt{loss}, \texttt{loss\_names} and \texttt{loss\_weights}; keep the
name \texttt{\textquotesingle{}crossentropy\textquotesingle{}} for the
classification loss or class balancing breaks \\
the four steps & unpack every output, compute every loss, combine with
the incoming \texttt{loss\_weights}, log one
\texttt{*\_\textless{}name\textgreater{}\_loss} per task \\
\texttt{evaluate()} & \textbf{override it inside \texttt{HW\_Net}} ---
\texttt{base.py} only unpacks multi-output models named
\texttt{MixNet}/\texttt{MIN2Net}; every other model hits
\texttt{np.argmax(test\_pred,\ axis=1)}, which is garbage for a list of
outputs \\
\end{longtable}
}

Do \textbf{not} rename your model to \texttt{MixNet*} to sneak into that
branch, and do not edit \texttt{base.py}. Overriding \texttt{evaluate()}
in \texttt{HW\_Net.py} is allowed --- it is one of your two files.

\texttt{mixnet/models/MixNet.py} is the reference implementation ---
read its steps side by side with the single-task ones in your file.

\subsection{Match the shapes}\label{match-the-shapes}

Nothing derives the shapes for you --- you set them yourself in
\texttt{experiments/configs/HW\_Net.py}, and the two must agree:

\begin{itemize}
\tightlist
\item
  \texttt{data\_params.data\_format} --- how \texttt{DataLoader}
  reshapes what Task 1 saved
\item
  \texttt{model\_params.input\_shape} --- what your \texttt{build()}
  receives, i.e.~the reshaped tensor \textbf{without} the leading trial
  axis
\end{itemize}

See \texttt{mixnet.utils.DataLoader.\_change\_data\_format} for details.

\subsection{Run it}\label{run-it-1}

\begin{Shaded}
\begin{Highlighting}[]
\BuiltInTok{cd}\NormalTok{ experiments}
\ExtensionTok{python}\NormalTok{ run\_HW\_Net.py}
\end{Highlighting}
\end{Shaded}

Results land in
\texttt{experiments/logs/HW\_Net/subject\_dependent\_2\_classes\_BCIC2a/}:

{\def\LTcaptype{none} % do not increment counter
\begin{longtable}[]{@{}
  >{\raggedright\arraybackslash}p{(\linewidth - 2\tabcolsep) * \real{0.5000}}
  >{\raggedright\arraybackslash}p{(\linewidth - 2\tabcolsep) * \real{0.5000}}@{}}
\toprule\noalign{}
\begin{minipage}[b]{\linewidth}\raggedright
File
\end{minipage} & \begin{minipage}[b]{\linewidth}\raggedright
Contents
\end{minipage} \\
\midrule\noalign{}
\endhead
\bottomrule\noalign{}
\endlastfoot
\texttt{S001\_all\_results.csv} & one row per fold ---
\texttt{test\_acc}, \texttt{f1-score}, timing, memory \\
\texttt{S001\_prediction\_results.npy} & \texttt{y\_true} /
\texttt{y\_pred} per fold \\
\texttt{S001\_fold01\_out\_weights.h5} & best checkpoint of that fold \\
\texttt{S001\_fold01\_out\_log.csv} & per-epoch training curve \\
\end{longtable}
}

\subsection{Tuning}\label{tuning}

Tune in \texttt{experiments/configs/HW\_Net.py} (\texttt{lr},
\texttt{batch\_size}, \texttt{dropout\_rate}, \texttt{es\_patience},
\ldots). Every key of \texttt{model\_params} is forwarded to your model
and re-applied as \texttt{self.\textless{}key\textgreater{}}, so you can
sweep hyper-parameters without touching model code.

\textbf{Select on validation, never on test.} Change the
\texttt{log\_path} suffix for each configuration so runs do not
overwrite each other --- \texttt{benchmark.py} discovers every directory
under \texttt{logs/} automatically.



\section{8. Task 3 --- benchmark against
MixNet}\label{task-3-benchmark-against-mixnet}

\begin{Shaded}
\begin{Highlighting}[]
\BuiltInTok{cd}\NormalTok{ experiments}
\ExtensionTok{python}\NormalTok{ benchmark.py}
\end{Highlighting}
\end{Shaded}

With no arguments it finds every run under \texttt{logs/}, averages the
5 folds within each subject, reports mean ± std across the 9 subjects,
and runs a paired comparison between a \texttt{MixNet} run and an
\texttt{HW\_Net} run. To be explicit:

\begin{Shaded}
\begin{Highlighting}[]
\ExtensionTok{python}\NormalTok{ benchmark.py }\AttributeTok{{-}{-}baseline}\NormalTok{ MixNet }\AttributeTok{{-}{-}candidate}\NormalTok{ HW\_Net }\AttributeTok{{-}{-}metric}\NormalTok{ test\_acc}
\end{Highlighting}
\end{Shaded}

It writes \texttt{benchmark\_per\_subject.csv} and
\texttt{benchmark\_summary.csv} and prints markdown tables you can paste
straight into your report.

\subsection{How to read the output}\label{how-to-read-the-output}

\begin{itemize}
\tightlist
\item
  \textbf{Per-subject table} --- the honest picture. BCIC2a has large
  between-subject variance; a method can win on average while losing on
  4 of 9 subjects.
\item
  \textbf{Paired t-test / Wilcoxon} --- is the difference more than
  noise? With only 9 subjects the power is low, so \emph{never} report a
  p-value without the per-subject table beside it.
\item
  \textbf{\texttt{candidate\ wins\ on\ k/9\ subjects}} --- often more
  informative than the mean.
\end{itemize}

\subsection{You are not required to beat
MixNet}\label{you-are-not-required-to-beat-mixnet}

MixNet is a published method with tuned hyper-parameters. A careful,
honest, well-analysed comparison that loses scores better than a win you
cannot explain or that came from a leak. What is graded is the quality
of the experiment and the analysis.



\section{9. Deliverables}\label{deliverables}

Push to your homework branch and submit its URL.

\begin{enumerate}
\def\labelenumi{\arabic{enumi}.}
\tightlist
\item
  \textbf{Code} --- your \texttt{HW\_prep.py}, \texttt{HW\_Net.py},
  \texttt{configs/HW\_Net.py}
\item
  \textbf{Results} --- \texttt{benchmark\_per\_subject.csv},
  \texttt{benchmark\_summary.csv}, and the \texttt{S*\_all\_results.csv}
  files for your runs (these are git-ignored by default; copy them into
  a \texttt{results/} folder at the repo root and commit that)
\item
  \textbf{\texttt{REPORT}} (3--5 A4 pages) at the repo root:

  \begin{itemize}
  \tightlist
  \item
    \textbf{Preprocessing} --- what you built and \emph{why}; how you
    prevented leakage
  \item
    \textbf{Model} --- architecture diagram or layer table, parameter
    count, design rationale
  \item
    \textbf{Results} --- the per-subject and summary tables, plus the
    paired test
  \item
    \textbf{Analysis} --- where you beat MixNet, where you lose, and
    your hypothesis why; which subjects were hardest and what they have
    in common
  \item
    \textbf{Reproduction} --- the exact commands you ran
  \end{itemize}
\end{enumerate}



\section{10. Rules}\label{rules}

\begin{enumerate}
\def\labelenumi{\arabic{enumi}.}
\tightlist
\item
  Do not modify anything under \texttt{mixnet/} other than your two
  files.
\item
  Do not train on session \texttt{E}, and do not select hyper-parameters
  on it.
\item
  Keep \texttt{k\_folds=5} and all 9 subjects, so the comparison stays
  paired.
\item
  Fit every data-driven transform on the training fold only.
\item
  Cite any external code or paper you draw on.
\item
  Report what you actually ran. Reproducibility is more important than
  your model's performance!
\end{enumerate}



\section{11. Grading}\label{grading}

{\def\LTcaptype{none} % do not increment counter
\begin{longtable}[]{@{}
  >{\raggedright\arraybackslash}p{(\linewidth - 2\tabcolsep) * \real{0.5000}}
  >{\raggedright\arraybackslash}p{(\linewidth - 2\tabcolsep) * \real{0.5000}}@{}}
\toprule\noalign{}
\begin{minipage}[b]{\linewidth}\raggedright
Criterion
\end{minipage} & \begin{minipage}[b]{\linewidth}\raggedright
Weight
\end{minipage} \\
\midrule\noalign{}
\endhead
\bottomrule\noalign{}
\endlastfoot
Preprocessing --- correct, leak-free, justified & 25\% \\
Model --- valid contract, sound design, justified & 25\% \\
Benchmark --- correct protocol, paired statistics, complete tables &
20\% \\
Analysis --- depth of interpretation & 20\% \\
Reproducibility --- the grader can re-run your commands and get your
numbers & 10\% \\
\end{longtable}
}



\section{12. Suggested schedule}\label{suggested-schedule}

{\def\LTcaptype{none} % do not increment counter
\begin{longtable}[]{@{}ll@{}}
\toprule\noalign{}
Week & Goal \\
\midrule\noalign{}
\endhead
\bottomrule\noalign{}
\endlastfoot
1 & Setup, data download, MixNet baseline reproduced and recorded \\
1-2 & Task 1 --- preprocessing running for all 9 subjects \\
1-2 & Task 2 --- model training end to end, first tuning pass \\
3 & Task 3 --- benchmark, \texttt{REPORT} \\
\end{longtable}
}

Good luck --- and start the baseline run early.

\end{document}
